\documentclass[../../main.tex]{subfiles}
\begin{document}

\begin{center}
  {\fontsize{48pt}{58pt}\selectfont \textbf{1981}\normalsize\textbf{年度}}
\end{center}

\begin{table}[htb]
  \centering
  \caption{出題テーマと難易度}
  \begin{tabular}{|c|p{6.2cm}|c|c|c|}
    \hline
    \textbf{問題番号} & \textbf{テーマ} & \textbf{思考量} & \textbf{計算量} & \textbf{総合} \\
    \hline
    第1問 & 放物線の弦の長さの最小化 & 普 & 難 & 普 \\
    \hline
    第2問 & 四面体の4辺上の点が作る平行四辺形 & 普 & 易 & 普 \\
    \hline
    第3問 & べき関数の不等式から定数を決定 & 普 & 普 & 普 \\
    \hline
    第4問 & 辺の長さを固定した四面体の最大体積 & 普 & 普 & 普 \\
    \hline
    第5問 & 回数制限のあるおみくじの期待回数 & 普 & 普 & 普 \\
    \hline
    第6問 & 指数関数の積分の不等式 & 難 & 普 & 難 \\
    \hline
  \end{tabular}
\end{table}

\begin{table}[htb]
  \centering
  \caption{解答}
  \renewcommand{\arraystretch}{1.4}
  \begin{tabular}{|c|c|p{9cm}|}
    \hline
    \textbf{問題番号} & \textbf{小問} & \textbf{答え} \\
    \hline
    第1問 & -- & 傾き$1$のとき最小長$8$ \\
    \hline
    第2問 & (1) & $k=l=m=n$ \\
    \cline{2-3}
    & (2) & 平行四辺形の対角線の交点は$OB,AC$の中点を結ぶ線分上にある \\
    \hline
    第3問 & -- & $k=\dfrac{m}{m+1}$ \\
    \hline
    第4問 & (1) & $\dfrac{\sqrt{3}}{27}a^3$ \\
    \cline{2-3}
    & (2) & $\dfrac{2\sqrt{3}}{27}a^3$ \\
    \hline
    第5問 & -- & $E=\dfrac{1-(1-p)^N}{p}$．$p=\dfrac{1}{5},N=10$なら$E\fallingdotseq4.46$ \\
    \hline
    第6問 & (1) & 指定された積分不等式は成立する（証明問題） \\
    \cline{2-3}
    & (2) & 指定された積分は$e^{3/2}$より小さい（証明問題） \\
    \hline
  \end{tabular}
\end{table}

\paragraph{試験形態}
6問．試験時間は要確認．

\end{document}
